\section{What the Formalisation Found}\label{sec:errata}

The HHL paper is a short letter with a long appendix, and its main claims survive mechanisation; they
are stated with $O(\cdot)$ and hold with other constants. Its details do not all survive.
Table~\ref{tab:errata} lists the statements of the arXiv version~\cite{hhl2008arxiv} that are false as
written, each with the correction and the formal theorem that refutes the claim or proves the
correction. Most sit in the appendix on optimality, which reduces quantum circuits to matrix
inversion; each optimality theorem is formalised up to the complexity-theoretic fact it cites.

The most instructive is Lemma~3, the Lipschitz estimate behind the paper's Theorem~1. The filter
functions mark an eigenvalue $\lambda$ as well-conditioned ($f$), ill-conditioned ($g$) or in between:
$f(\lambda)=1/(2\kappa\lambda)$ and $g(\lambda)=0$ for $\lambda\ge1/\kappa$, $f=0$ and $g=\frac12$ below
$1/(2\kappa)$, and on $[\frac1{2\kappa},\frac1\kappa)$ they turn through a quarter circle of radius
$\frac12$. Lemma~3 claims
\[
|f(\lambda)-f(\tilde\lambda)|^2+|g(\lambda)-g(\tilde\lambda)|^2\;\le\;c\,\kappa^2(\lambda-\tilde\lambda)^2\bigl(|f(\lambda)|^2+|g(\lambda)|^2\bigr)
\]
with $c=\pi^2/2$. At $\lambda=\frac1{2\kappa}$, $\tilde\lambda=\frac1\kappa$ the left side is $\frac14+\frac14=\frac12$,
while $\kappa^2(\lambda-\tilde\lambda)^2=\frac14$ and $f(\lambda)^2+g(\lambda)^2=\frac14$, so the right side is
$c/16=\pi^2/32<\frac12$, for every~$\kappa$. The constant $c=\pi^2$
holds, and no smaller constant does (Listing~\ref{lst:lemma3}). The paper uses the lemma only inside an
$O(\kappa/t)$ bound, so its Theorem~1 is unaffected; with the corrected constant we prove it with
the explicit constant $12\pi^2\sqrt{4+2\pi^2}$.

%%LST Algorithms/HHL_Filter.thy lemma3_paper_constant lemma3 lemma3_sharp
\begin{lstlisting}[float=tp,caption={Lemma~3 of the HHL paper: the counterexample to the paper's constant, the corrected constant, and its sharpness (for fixed \isa{$\kappa$ > 0}; \isa{filt\_f} and \isa{filt\_g} are the filter functions)},label={lst:lemma3}]
theorem lemma3_paper_constant:
  "pi~$^{2}$~ / 2 * ~$\kappa^{2}$~ * (1 / (2 * ~$\kappa$~) - 1 / ~$\kappa$~)~$^{2}$~ * ((filt_f ~$\kappa$~ (1 / (2 *
      ~$\kappa$~)))~$^{2}$~ + (filt_g ~$\kappa$~ (1 / (2 * ~$\kappa$~)))~$^{2}$~) < (filt_f ~$\kappa$~ (1 / (2 *
      ~$\kappa$~)) - filt_f ~$\kappa$~ (1 / ~$\kappa$~))~$^{2}$~ + (filt_g ~$\kappa$~ (1 / (2 * ~$\kappa$~)) - filt_g
      ~$\kappa$~ (1 / ~$\kappa$~))~$^{2}$~"
theorem lemma3:
  "(filt_f ~$\kappa$~ a - filt_f ~$\kappa$~ b)~$^{2}$~ + (filt_g ~$\kappa$~ a - filt_g ~$\kappa$~ b)~$^{2}$~ ~$\le$~ pi~$^{2}$~
      * ~$\kappa^{2}$~ * (a - b)~$^{2}$~ * ((filt_f ~$\kappa$~ a)~$^{2}$~ + (filt_g ~$\kappa$~ a)~$^{2}$~)"
theorem lemma3_sharp:
  assumes c: "c < pi~$^{2}$~"
  shows "~$\exists$~a b. c * ~$\kappa^{2}$~ * (a - b)~$^{2}$~ * ((filt_f ~$\kappa$~ a)~$^{2}$~ + (filt_g ~$\kappa$~ a)~$^{2}$~)
      < (filt_f ~$\kappa$~ a - filt_f ~$\kappa$~ b)~$^{2}$~ + (filt_g ~$\kappa$~ a - filt_g ~$\kappa$~ b)~$^{2}$~"
\end{lstlisting}

\begin{table}[t]
\caption{Errata in the HHL paper (arXiv:0811.3171v3). ``Formal'' names the theorem that refutes the claim
or proves the correction.}\label{tab:errata}
\centering\scriptsize
\renewcommand{\arraystretch}{1.12}
\newcolumntype{L}[1]{>{\raggedright\arraybackslash}p{#1}}
\begin{tabular}{@{}L{1.3cm}L{2.6cm}L{4.5cm}L{2.75cm}@{}}
\toprule
Where & Claim & Correction & Formal\\
\midrule
Lemma 3 & $c=\pi^2/2$ & false for every $\kappa$; $c=\pi^2$ holds and is sharp & \isa{lemma3\_paper\_constant}, \isa{lemma3}, \isa{lemma3\_sharp}\\
Proof of Thm.~4 & $m=\frac2\delta\log(2n)/\log\log n$ rounds suffice & too few: for $n_0=2^{64}$, $\delta=\frac12$ the bound on $T_m$ still exceeds $n_0$; $m\ge\frac2\delta\ln(2n_0)$ suffices & \isa{paper\_rounds\_insufficient}, \isa{corrected\_rounds}\\
(A34) & the dilation satisfies $\kappa^{-1}I\le A\le I$ & impossible: its diagonal is zero & \isa{herm\_embed\_diag}\\
(A37) & the first qubit of $\vec y$ decides acceptance & the top half of $\vec y$ is always $0$; permuting the dilation fixes it & \isa{first\_half\_perm}, \isa{red\_first\_qubit}\\
Thm.~6(2) & precision $\varepsilon=2^{-n-2}$ & cannot separate counts $C$ and $C+1$ once the clock factor ($\le\frac14$) is included & \isa{paper\_precision\_insufficient}\\
App.~A5 & $1/(3\varepsilon^2)$ samples & too few: $\varepsilon=\frac14$, $p=\frac12$ and $6$ samples fail with probability $\frac{11}{16}$; $3/\varepsilon^2$ suffice & \isa{paper\_sample\_count\_too\_small}, \isa{estimate\_expectation}\\
\S IV & preconditioned system $(AB)c=Bb$ & neither left nor right preconditioning: in a $2\times2$ example its solution gives no solution of $Ax=b$ & \isa{paper\_precondition\_wrong}\\
App.~A5 & undo gates $U^\dagger_{3T+1-t}$ & must be $U^\dagger_{T+1-t}$ & construction \isa{red\_U}, \isa{red\_inverse}\\
App.~A5 & $\kappa\le2T$ & we prove $\kappa\le2T+1$ & \isa{red\_A\_bounds}\\
\bottomrule
\end{tabular}
\end{table}

The first seven rows are refuted by explicit counterexamples; the last two are corrections built into
the formal construction. Three observations generalise. First, none of the errors invalidates a main
claim. Second, the errors sit where a proof sketch compresses a computation: an evaluation of
$f^2+g^2$ at an endpoint, an index, a number of rounds, a count of samples. A reader checks the
argument; a proof assistant checks the arithmetic, and the arithmetic is where the sketch was wrong.
Third, a formal proof settles disagreements between readers. A recent re-derivation of the error
analysis~\cite{liphillips2024} remarks that the factor $\pi/2$ of the paper's Lemma~2 ``is not correct
for certain values of~$\kappa$''; with the paper's choice $\kappa'=2\kappa$, our theorem
\isa{hflag\_lipschitz} shows that $\lambda\mapsto\ket{h(\lambda)}$ is $\frac\pi2\kappa$-Lipschitz for
every $\kappa>0$, so Lemma~2 is correct as stated.

All of these were found during mechanisation, in the course of proving the paper's statements rather
than by a separate review. Every row of the table has a formal witness. One further doubt we could only argue
informally: the paper's bound $\kappa\le2T$ appears to fail when $3T$ is even, where
$\kappa=\coth\frac1{2T}>2T$. There we prove the weaker $\kappa\le2T+1$ and do not count the claim as refuted.
